\documentclass[11pt,a4paper]{article}

\usepackage[utf8]{inputenc}
\usepackage[T1]{fontenc}
\usepackage[round,sort,comma]{natbib}
\usepackage{amsmath,amssymb}
\usepackage{booktabs}
\usepackage{listings}
\usepackage{xcolor}
\usepackage[hidelinks,backref=section]{hyperref}
\usepackage{microtype}
\usepackage{geometry}
\geometry{margin=1in}

\definecolor{skyblue}{RGB}{48, 141, 222}
\definecolor{northblue}{RGB}{49, 126, 204}
\definecolor{darkblue}{RGB}{1, 81, 144}
\definecolor{fade_darkblue}{RGB}{179, 222, 255}
\definecolor{arizona}{RGB}{83, 76, 76}
\definecolor{emerald}{RGB}{0, 129, 133}
\definecolor{forest}{RGB}{12, 153, 69}
\definecolor{fade_emerald}{RGB}{161, 227, 230}
\definecolor{crayon}{RGB}{255, 207, 0}
\definecolor{peel}{RGB}{255, 102, 0}
\definecolor{fade_peel}{RGB}{255, 209, 179}
\definecolor{chocolate}{RGB}{117, 48, 3}
\definecolor{prisma}{RGB}{64, 59, 121}
\definecolor{darkpurple}{RGB}{104, 59, 121}
\definecolor{lush}{RGB}{2, 163, 68}
\definecolor{fade_lush}{RGB}{166, 237, 193}
\definecolor{flame}{RGB}{232, 51, 81}

\hypersetup{
    colorlinks = true,
    linkcolor = skyblue,
    citecolor = chocolate,
    filecolor = peel,
    urlcolor = flame
}

\renewcommand*{\backref}[1]{}
\renewcommand*{\backrefalt}[4]{%
  \ifcase #1 %
    {\hfill \it}%
  \or
    {\hfill {\it In sec. {\hypersetup{linkcolor=chocolate}#2}.}}%
  \else
    {\hfill {\it In sec. {\hypersetup{linkcolor=chocolate}#2}.}}%
  \fi
}

\lstset{
  basicstyle=\ttfamily\small,
  frame=single,
  breaklines=true,
  keepspaces=true,
}

\setlength{\parindent}{0pt}
\setlength{\parskip}{0.5em}

\title{Program Generation via Reinforcement Learning\\with Enumeration Warm-Start}
\author{Yiding Song}
\date{\today}

\begin{document}
\maketitle

\section*{Problem Statement}

We are interested in the following problem. Suppose we have access to:
\begin{itemize}
  \item A set $F$ of functions $f_1, f_2, \ldots, f_n$ with types
        $\tau_1 \to \upsilon_1, \tau_2 \to \upsilon_2, \ldots, \tau_n \to \upsilon_n$,
        where we can query the function implementation (i.e.\ $f_i(x)$ can be computed
        for arbitrary well-typed inputs $x$) but not necessarily the function description
        (e.g.\ we do not know the mathematical or programmatic form of the functions).
  \item A reward function $R: \mathrm{program} \to \mathbb{R}$ that maps a program to a
        real-valued reward, where a program is a well-typed composition of functions from
        $F$ (e.g.\ this can reward a function for desirable properties such as being
        non-constant).
\end{itemize}

{\it How do we synthesise a large number of programs that are diverse and achieve high rewards?}

This document describes a two-phase program synthesis solution to the problem: bottom-up enumeration
followed by reinforcement learning.

\section*{Overview}

The pipeline has two phases:

\begin{enumerate}
  \item \textbf{Bottom-up enumeration} --- exhaust all programs up to a certain size,
        quotienting out semantic duplicates (as measured by function behaviour on a test
        suite), creating a corpus of behaviourally diverse programs.
  \item \textbf{RL exploration} --- a neural policy learns from the corpus (behavioural
        cloning warm-start), then generates new programs via policy gradient, using a
        priority queue buffer of the best programs found so far
        \citep{abolafia2018priority}.
\end{enumerate}

\section{The Program Grammar}\label{sec:grammar}

Programs are expressions in a simply-typed lambda calculus extended with conditionals and
a fixed set of primitive functions. The expression language is:
\[
  e \;::=\; c \;\mid\; x \;\mid\; f(e_1, \ldots, e_k)
         \;\mid\; \lambda(x_1, \ldots, x_k).\, e
         \;\mid\; \mathrm{if}\ e_1\ \mathrm{then}\ e_2\ \mathrm{else}\ e_3
\]
where $c$ ranges over literals (integers, booleans, empty lists), $x$ over variables, and
$f$ over the primitive functions in $F$.

\subsection{Types}\label{sec:types}

The base types are $\mathrm{int}$ and $\mathrm{bool}$. From these, list types
$\mathrm{list}[\sigma]$ and function types $\sigma \to \tau$ are formed. In the current
instantiation the concrete types in use are:
\[
  \mathrm{int},\quad \mathrm{bool},\quad \mathrm{list}[\mathrm{int}],\quad
  \mathrm{list}[\mathrm{bool}],\quad \mathrm{list}[\mathrm{list}[\mathrm{int}]]
\]
All programs synthesised by the pipeline have type
$\mathrm{list}[\mathrm{int}] \to \mathrm{list}[\mathrm{int}]$.

\subsection{Polymorphic Functions}\label{sec:polymorphic}

Some primitives in $F$ are \textbf{polymorphic}: their types contain type variables that
can be instantiated to any concrete type. For example:
\[
  \mathrm{map} : (T_1 \to T_2) \to \mathrm{list}[T_1] \to \mathrm{list}[T_2]
\]
Here $T_1$ and $T_2$ are type variables. A \textbf{type substitution} $\theta$ maps each
type variable to a concrete type; applying $\theta$ to the signature of $\mathrm{map}$
yields a monomorphic instantiation. For instance,
$\theta = \{T_1 \mapsto \mathrm{int},\, T_2 \mapsto \mathrm{bool}\}$ gives:
\[
  \mathrm{map}_\theta : (\mathrm{int} \to \mathrm{bool})
    \to \mathrm{list}[\mathrm{int}] \to \mathrm{list}[\mathrm{bool}]
\]
The set of valid substitutions for a function is determined by the concrete types in use.
Each instantiation is treated as a separate monomorphic function during enumeration.

\subsection{Higher-Order Functions}\label{sec:hof}

A function is \textbf{higher-order} if one or more of its arguments has a function type.
\texttt{map} is higher-order: its first argument is a function $T_1 \to T_2$. Such
arguments cannot be looked up directly in the program bank (which stores complete
expressions, not open function values); instead they must be synthesised as lambda
expressions, as described in the Higher-Order Applications section below.

\section{Part 1: Bottom-Up Enumeration}\label{sec:enum}

\subsection{Size Metric}\label{sec:size}

Every AST node $e$ is assigned an integer size $\lvert e \rvert$ inductively:

\begin{center}
\begin{tabular}{ll}
\toprule
Node & Size \\
\midrule
Literal, variable, empty list
  & $1$ \\
Application $(f\ a_1\ \cdots\ a_k)$
  & $1 + \lvert a_1 \rvert + \cdots + \lvert a_k \rvert$ \\
Lambda $(\lambda\ (p_1\ \cdots\ p_k)\ e)$
  & $1 + \lvert e \rvert$ \\
Conditional $(\mathrm{if}\ c\ t\ e)$
  & $1 + \lvert c \rvert + \lvert t \rvert + \lvert e \rvert$ \\
\bottomrule
\end{tabular}
\end{center}

Enumeration proceeds by increasing size: programs of size $s$ are built exclusively from
sub-expressions of size less than $s$.

\subsection{Observational Equivalence and the Program Bank}\label{sec:oe}

We remove program duplicates by comparing their behaviour on a test suite. Fix a test suite $T = \{t_1, \ldots, t_m\}$ of $m$
concrete inputs. The \textit{fingerprint} of a closed program $p$ of type
$\tau \to \upsilon$ is:
\[
  \Phi(p) = \bigl(p(t_1),\ \ldots,\ p(t_m)\bigr) \in (\upsilon \cup \{\dagger\})^m
\]
where $\dagger$ records a failed evaluation (runtime error or type error). Two programs
$p, q : \tau \to \upsilon$ are observationally equivalent if $\Phi(p) = \Phi(q)$.

\paragraph{Open terms and lambda arguments.}
However, not all programs are closed. To enumerate an expression like
$\mathrm{map}\ (\lambda y.\ e)\ x$, the enumerator must synthesise candidate bodies $e$.
These bodies are \emph{open terms}: they may reference $y$ (the lambda parameter) freely,
alongside the outer $x$. Storing them in the bank allows the same body to be reused
across different higher-order functions, avoiding redundant work.

In general, for an open term $e$ with free variables
$y_1 : \sigma_1, \ldots, y_k : \sigma_k$ (beyond the primary input $x$), the fingerprint
is computed over $T \times V_1 \times \cdots \times V_k$, where each $V_i$ is a finite
set of probe values for type $\sigma_i$:
\[
  \Phi(e)
    = \bigl(e(t,\, v_1,\ldots,v_k)\bigr)_{(t,\, v_1,\ldots,v_k)
        \,\in\, T \times V_1 \times \cdots \times V_k}
\]
Two open terms are considered observationally equivalent iff they agree on this full
Cartesian product.

The \textbf{program bank} $\mathcal{B}$ maps each pair $(\tau, s)$ to a set of programs,
containing exactly one representative per observational equivalence class:
\[
  \mathcal{B}[\tau][s]
    = \bigl\{[p]_\sim \;\big|\; \vdash p : \tau,\; \lvert p \rvert = s\bigr\}
\]

\paragraph{Default test suite} ($m = 10$):
\begin{lstlisting}
[], [0], [3,1,2], [1,1,1,1], [5,4,3,2,1],
[1,2,3,4,5,6,7,8], [10,-3,7,7,0], [2,8,3,8,2,3],
[0,1,0,1,0], [42]
\end{lstlisting}

\subsection{Base Case}\label{sec:base}

The bank is seeded with size-1 expressions:
\begin{itemize}
  \item \textbf{Integer hole} $\square_{\mathbb{Z}}$ of type $\mathrm{int}$ (see Integer Holes below)
  \item \textbf{Boolean constants} $\top, \bot$ of type $\mathrm{bool}$
  \item \textbf{Empty lists} $[\,]$ of type $\mathrm{list}[\sigma]$ for each list type $\sigma$
  \item \textbf{Input variable} $x : \mathrm{list}[\mathrm{int}]$ (since all of our
        programs are of type $\mathrm{list}[\mathrm{int}] \to \mathrm{list}[\mathrm{int}]$,
        this is the variable that will be bound to the input list in a program like
        \texttt{length x})
\end{itemize}

\subsection{Integer Holes}\label{sec:holes}

Instead of seeding a fixed set of integer literals $C_{\text{seed}} = \{0,\ldots,10\}$,
the base case includes a single \textbf{integer hole} $\square_{\mathbb{Z}}$ of type
$\mathrm{int}$. A hole is an uninstantiated integer position; its concrete value is
determined by a \textbf{substitution} $\sigma$, an ordered list of integers indexed by
pre-order traversal position in the AST. Therefore, concrete programs are a tuple $(\pi, \sigma)$ where $\pi$ is a valid program sketch and $\sigma$ is the substitution of concrete values for the integer holes. This allows us to enumerate / diversify programs involving integer constants much more efficiently.

\paragraph{Fingerprinting.}
To compute the fingerprint of a sketch $\pi$ containing $k$ holes, the compiler
substitutes the $i$-th hole (in pre-order) with
$\mathrm{PROBE}[i \bmod |\mathrm{PROBE}|]$ where we arbitrarily chose the random sequence
\[
  \mathrm{PROBE} = [5, 3, 7, 2, 9, 0, 8, 1, 6, 4].
\]
The resulting fingerprint has length $m$ (same as concrete programs). Two sketches
that produce identical outputs under the probe substitution are considered
observationally equivalent and deduplicated. False equivalences are possible but
rare given the diversity of the test suite.

\paragraph{Corpus expansion.}
After enumeration, a sketch $\pi$ can be expanded into concrete programs by
sampling $N_{\mathrm{fill}}$ substitutions
$\sigma \sim \mathrm{Uniform}(C_{\mathrm{seed}}^k)$ and retaining those that pass
the quality filter. This expansion is a standalone utility; the RL pipeline works
directly on sketches.

\subsection{First-Order Applications}\label{sec:first-order}

For a first-order function $f : A_1 \times \cdots \times A_k \to R$ and target size $s$,
the arguments must satisfy $\sum_{i=1}^k s_i = s - 1$ with each $s_i \geq 1$. For every
such composition:
\[
  \mathcal{B}[R][s]
    \;\mathrel{+}=\;
    \bigl\{f(a_1, \ldots, a_k) \;\big|\; a_i \in \mathcal{B}[A_i][s_i]\bigr\}
    \!\Big/\!\sim
\]
Polymorphic functions are pre-expanded into all valid monomorphic instantiations (as
described in the Grammar section); each instantiation is then enumerated independently
using the formula above.

\subsection{Higher-Order Applications and Nested Lambdas}\label{sec:ho-apps}

When an argument position requires a function type $A_1 \times \cdots \times A_k \to B$,
the lambda argument is synthesised inline rather than looked up in the flat bank. Given a size
budget $b$ for the lambda, an extended typing context
$\Gamma' = \Gamma \cup \{y_i : A_i\}_{i=1}^k$ is formed, and a child bank $\mathcal{B}'$
is enumerated recursively over $\Gamma'$ up to size $b - 1$. The lambda candidates are
then:
\[
  \bigl\{(\lambda\ (y_1 \cdots y_k)\ e) \;\big|\; e \in \mathcal{B}'[B][b-1]\bigr\}
\]
Child banks are keyed by $(\Gamma', b)$ and cached so that the same lambda type appearing
as an argument to several functions at the same size is only built once.

\paragraph{Nesting depth.}
$\ell$ counts the number of enclosing lambda scopes. Higher-order calls within a child
context are only attempted when $\ell \leq \ell_{\max}$ (default $\ell_{\max} = 2$),
preventing combinatorial explosion from arbitrarily deep lambda nesting while still
enabling patterns with three levels of nesting such as
$\mathrm{map}\ (\lambda y.\ \mathrm{map}\ (\lambda z.\ \mathrm{map}\ (\lambda w.\ \cdots)\ z)\ y)\ x$.

\subsection{Quality Filtering}\label{sec:quality}

After enumeration, the corpus is obtained by retaining programs that pass three
conditions. Let $\mathcal{F}(p) \subseteq \{1,\ldots,m\}$ denote the indices on which
$p$ does not crash. A program $p$ is retained iff:
\begin{enumerate}
  \item $|\mathcal{F}(p)| \geq 3$ \quad (non-crashing on sufficiently many inputs)
  \item $|\{\,p(t_i) \mid i \in \mathcal{F}(p)\}| \geq 2$ \quad (non-constant output)
  \item $\dfrac{|\{\,p(t_i) \mid i \in \mathcal{F}(p)\}|}{|\mathcal{F}(p)|}
        \geq \nu_{\min}$ \quad (output variability, default $\nu_{\min} = 0.3$)
\end{enumerate}

\subsection{Key Parameters}\label{sec:params}

\begin{center}
\begin{tabular}{lll}
\toprule
Parameter & Default & Role \\
\midrule
$s_{\max}$            & 8         & Maximum AST size enumerated \\
$\ell_{\max}$         & 2         & Maximum lambda nesting depth for HOF arguments \\
$N_{\mathrm{fill}}$   & 10        & Random substitutions sampled per sketch during corpus expansion \\
$\nu_{\min}$          & 0.3       & Corpus variability threshold \\
$T$                   & 10 inputs & Test suite for fingerprinting and deduplication \\
\bottomrule
\end{tabular}
\end{center}

\section{Part 2: Reinforcement Learning}\label{sec:rl}

The RL phase trains a neural policy on the enumerated corpus to generate good programs
recursively from the root, making one decision per AST node. This is framed as a Markov
Decision Process (MDP) over a grammar-constrained action space.

\subsection{MDP Formulation}\label{sec:mdp}

\paragraph{State.}
A synthesis state $s$ is a tuple $(\tau, \Gamma, f_{\text{par}}, i, \delta, \ell)$:

\begin{center}
\begin{tabular}{ll}
\toprule
Component & Description \\
\midrule
$\tau$                  & Required type at the current node \\
$\Gamma$                & Typing context: variables in scope with their types \\
$f_{\text{par}},\ i$   & Parent function and argument index (contextual features) \\
$\delta$                & Remaining depth budget (decremented at each recursive call) \\
$\ell$                  & Current lambda nesting depth \\
\bottomrule
\end{tabular}
\end{center}

\paragraph{Actions.}
The valid action set $\mathcal{A}(s)$ depends on $s$:

\begin{center}
\begin{tabular}{ll}
\toprule
Action & Condition \\
\midrule
Integer hole $\square_{\mathbb{Z}}$ & $\tau = \mathrm{int}$ \\
Literal boolean      & $\tau = \mathrm{bool}$ \\
Empty list           & $\tau$ is a list type \\
Variable $y$         & $y \in \mathrm{dom}(\Gamma)$ with $\Gamma(y) = \tau$ \\
Apply $f$ (instance) & $\delta > 0$; return type of $f$ matches $\tau$;
                       if $f$ is higher-order, $\ell < \ell_{\max}^{\mathrm{RL}} = 2$ \\
Lambda               & $\tau$ is a function type \\
If                   & $\delta > 0$ \\
\bottomrule
\end{tabular}
\end{center}

\paragraph{Episode.}
Generation begins from the state
$s_0 = (\mathrm{list}[\mathrm{int}] \to \mathrm{list}[\mathrm{int}],\
\varnothing,\ \cdot,\ \cdot,\ \delta_{\max},\ 0)$
and proceeds by recursively expanding each chosen action into child states:
\begin{itemize}
  \item $\mathrm{Apply}\ f(A_1, \ldots, A_k) \to R$: recurse on
        $(\tau \leftarrow A_i,\ \delta - 1)$ for each $i$
  \item $\mathrm{Lambda}\ (\tau_1 \times \cdots \times \tau_k \to \sigma)$: recurse on
        $(\tau \leftarrow \sigma,\ \Gamma \cup \{y_i : \tau_i\},\ \delta - 1,\ \ell + 1)$
  \item $\mathrm{If}$: recurse on $(\mathrm{bool},\ \delta-1)$ then twice on
        $(\tau,\ \delta-1)$
\end{itemize}
If any recursive call encounters an empty action set, the episode fails and returns no
program.

\subsection{Reward}\label{sec:reward}

Given a completed program $p$, let $\Phi(p)$ denote its fingerprint on $T$, and let
$\mathcal{K}$ be the set of fingerprints already known (from the enumeration corpus and
prior RL discoveries). The reward is:
\[
  R(p) = \mathrm{var}(\Phi(p)) \times
    \begin{cases}
      1.0 & \text{if } \Phi(p) \notin \mathcal{K} \\
      0.1 & \text{otherwise}
    \end{cases}
\]
where $\mathrm{var}(\Phi(p)) = |\{p(t_i)\}| / m$ is the fraction of distinct outputs.
Programs with $|\mathcal{F}(p)| < 3$ or $|\{p(t_i)\}| < 2$ receive $R(p) = 0$ and are
discarded.

The novelty multiplier incentivises exploration of behaviours not yet covered, while the
$0.1$ floor ensures the policy is not penalised for rediscovering known programs.

\subsection{Priority Queue Buffer}\label{sec:pq}

The buffer follows the priority queue training scheme of \citet{abolafia2018priority}.
$\mathcal{B}_{\mathrm{buf}}$ maintains the top-$K$ programs by reward (default
$K = 5000$). It is a bounded min-heap: a new program $p$ is admitted only if
$R(p) > R_{\min}(\mathcal{B}_{\mathrm{buf}})$, evicting the current minimum. Duplicate
fingerprints are excluded.

Training samples are drawn uniformly from $\mathcal{B}_{\mathrm{buf}}$ with no recency or
priority weighting.

\subsection{Policy Network}\label{sec:policy}

The policy $\pi_\theta$ maps a state $s$ to a distribution over actions via a masked
softmax. A network $g_\theta$ produces a logit for each action; invalid actions are
excluded by setting their logits to $-\infty$ before normalising:
\[
  \pi_\theta(a \mid s)
    = \frac{\exp\bigl(g_\theta(s)_a\bigr)}
           {\displaystyle\sum_{a' \in \mathcal{A}(s)} \exp\bigl(g_\theta(s)_{a'}\bigr)},
  \qquad a \in \mathcal{A}(s)
\]
where $g_\theta : \mathcal{S} \to \mathbb{R}^{|\mathcal{A}|}$ is a two-layer feedforward
network:
\[
  g_\theta(s) = W_2\, \mathrm{ReLU}(W_1\, \mathrm{enc}(s) + b_1) + b_2
\]

\paragraph{State encoder.}
Six features of $s$ are each embedded or projected into $\mathbb{R}^d$, then concatenated
and projected back to $\mathbb{R}^d$:

\begin{center}
\begin{tabular}{ll}
\toprule
Feature & Encoding \\
\midrule
$\tau$              & Learned embedding over a type vocabulary \\
$f_{\mathrm{par}}$  & Learned embedding over the function vocabulary \\
$i$ (arg index)     & Learned embedding, clamped to $[0,7]$ \\
$\delta$            & Learned embedding, clamped to $[0,15]$ \\
$\ell$              & Learned embedding, clamped to $[0,3]$ \\
$\Gamma$            & Linear projection of a per-type variable-count vector \\
\bottomrule
\end{tabular}
\end{center}

Default $d = 64$, hidden dimension 128. Invalid actions are masked to $-\infty$ before
softmax; the valid set $\mathcal{A}(s)$ is cached by $(\tau, \Gamma, \delta > 0, \ell)$.

\subsection{Training Phase 1: Behavioural Cloning Warm-Start}\label{sec:warmstart}

Each corpus program $p$ induces a sequence of state--action pairs $\{(s_t, a_t)\}$ by a
DFS traversal of its AST, recording at each node the state that a top-down policy would
have been in and the action it would have taken. Aggregating over all corpus programs
gives a dataset $\mathcal{D}$. The policy is pre-trained by minimising:
\[
  \mathcal{L}_{\mathrm{BC}}(\theta)
    = -\mathbb{E}_{(s,\, a)\,\sim\,\mathcal{D}}\bigl[\log \pi_\theta(a \mid s)\bigr]
\]
(50 epochs, batch size 64, Adam at $\eta = 10^{-3}$). Transitions involving actions
outside the vocabulary are dropped.

\subsection{Training Phase 2: RL Loop}\label{sec:rl-loop}

Each iteration alternates between a sampling phase and a gradient phase.

\paragraph{Sampling} (32 episodes per iteration): run an episode under $\pi_\theta$
to obtain a sketch $\pi$ (which may contain integer holes $\square_{\mathbb{Z}}$);
sample $N_{\mathrm{fill}} = 10$ random substitutions
$\sigma \sim \mathrm{Uniform}(C_{\mathrm{seed}}^k)$ where $k$ is the number of holes
in $\pi$; select the $\sigma^*$ maximising $R$; compute $\Phi(\pi[\sigma^*])$ and
$R(\pi[\sigma^*])$; if $R > 0$, attempt insertion of the concrete program into
$\mathcal{B}_{\mathrm{buf}}$; if $\Phi$ is novel, add it to $\mathcal{K}$.

\paragraph{Gradient update} (8 steps per iteration): sample a batch from
$\mathcal{B}_{\mathrm{buf}}$ and minimise the reward-weighted log-likelihood loss
\citep{abolafia2018priority}:
\[
  \mathcal{L}_{\mathrm{RL}}(\theta)
    = -\mathbb{E}_{p\,\sim\,\mathcal{B}_{\mathrm{buf}}}
        \left[ R(p) \sum_{t} \log \pi_\theta(a_t \mid s_t) \right]
\]
No baseline is subtracted. Gradient norms are clipped to 1.0 before the Adam step
($\eta = 10^{-4}$).

\subsection{Trajectory Extraction}\label{sec:trajectory}

The dataset $\mathcal{D}$ (used in both BC and RL) is constructed by a DFS that is the
inverse of the episode procedure. Given a complete AST, each node yields a state--action
pair:
\begin{itemize}
  \item $\lambda$-node $\to$ action \textbf{Lambda}, recurse into body at $\ell + 1$
  \item Application $f(a_1, \ldots, a_k)$ $\to$ action \textbf{Apply} $f$, recurse into
        each $a_i$ with $f_{\mathrm{par}} = f$, $i$ set accordingly
  \item If-node $\to$ action \textbf{If}, recurse into condition then branches
  \item Leaf $\to$ the appropriate literal or variable action
\end{itemize}
For polymorphic functions, the type instantiation is recovered by matching the inferred
return type of $f$ against $\tau$ in the current state.

\section{Overall Pipeline}\label{sec:pipeline}

This section describes the procedure used to generate 5{,}699{,}938 unique and behaviourally diverse programs using the list manipulation DSL primitives provided in \citet{rule2024symbolic}, linking all of the above methods.

\begin{itemize}

  \item \textbf{Phase 1 --- Bottom-Up Enumeration}
    (cf.~\S\ref{sec:enum}).
    A \texttt{BottomUpEnumerator} is run with parameters $s_{\max} = 8$,
    $\ell_{\max} = 2$, and $\nu_{\min} = 0.3$ on the default 10-input test suite
    and seed constants $C_{\mathrm{seed}} = \{0, 1, \ldots, 10\}$.
    The enumerator fills the program bank $\mathcal{B}$ with all
    behaviourally distinct \emph{sketches} (ASTs that may contain integer
    holes $\square_{\mathbb{Z}}$) up to size $s_{\max}$, yielding
    643{,}837 sketches in total.
    Two sub-products are retained:
    \begin{itemize}
      \item \textbf{\texttt{quality\_corpus}} (581{,}984 sketches) --- those
            that pass the quality filter (\S\ref{sec:quality}): at least 3
            non-crashing inputs, at least 2 distinct outputs, and output
            variability $\geq \nu_{\min}$.  Used for warm-start training and
            buffer seeding.
      \item \textbf{\texttt{all\_sketches}} (643{,}837 sketches) --- the full
            bank regardless of quality, flattened across all types and sizes.
            Used as the input to sketch expansion in Phase~2.
    \end{itemize}
    Quality sketches are serialised to \texttt{enum\_sketches.json}; the full
    enumeration result is checkpointed to \texttt{enum\_corpus.pkl} so that
    Phases~1 and~2 can be skipped on subsequent runs.

  \item \textbf{Phase 2 --- Sketch Expansion}
    (cf.~\S\ref{sec:holes}).
    Every sketch in \texttt{all\_sketches} is concretised by the
    \texttt{expand\_sketches} utility.  For a sketch $\pi$ containing $k$
    holes, substitutions $\sigma \in C_{\mathrm{seed}}^k$ are enumerated
    exhaustively when $|C_{\mathrm{seed}}|^k$ is small, or otherwise sampled
    uniformly, up to a per-sketch budget derived from
    \[
      N_{\mathrm{sub}} = \left\lceil \frac{N_{\mathrm{target}}}{|\texttt{all\_sketches}|} \right\rceil,
      \qquad N_{\mathrm{target}} = 6{,}000{,}000.
    \]
    Each concrete program is fingerprinted (\S\ref{sec:oe}), checked against a
    running deduplication table, and retained only if it passes the quality
    filter.  In practice, most sketches contain zero integer holes and therefore
    contribute exactly one concrete program; the resulting deduplicated set
    \texttt{enum\_concrete} contains 1{,}793{,}723 programs, written to
    \texttt{enum\_corpus.json}.  All fingerprints are added to
    \texttt{corpus\_fingerprints}, which tracks every behaviour seen so far.

  \item \textbf{Phase 3 --- RL Collection and Expansion}
    (cf.~\S\ref{sec:rl}).
    \begin{itemize}
      \item \emph{Warm-start} (\S\ref{sec:warmstart}).  The policy
            network is initialised via behavioural cloning on
            \texttt{quality\_corpus}, subsampled to 100{,}000 programs
            (yielding 879{,}141 state--action transitions).  Training runs for
            5 epochs with batch size 128 and Adam at $\eta = 10^{-3}$, using
            GPU acceleration when available (MPS on Apple Silicon, CUDA
            otherwise).  The trained weights are saved to
            \texttt{policy\_warmstart.pt} so that warm-start can be skipped on
            subsequent runs.

      \item \emph{Buffer seeding} (\S\ref{sec:pq}).  The
            priority queue buffer (capacity $K = 50{,}000$) is seeded with
            programs from \texttt{quality\_corpus} whose trajectories can be
            extracted; each is inserted with reward computed against
            \texttt{corpus\_fingerprints}.  In this run, the buffer was filled
            to capacity (50{,}000 programs).

      \item \emph{RL loop} (\S\ref{sec:rl-loop}).  Up to
            $5{,}000{,}000$ iterations are run, stopping early once
            100{,}000 novel sketches have been collected.  Each iteration
            consists of:
            \begin{enumerate}
              \item \emph{Sampling} (32 episodes): an \texttt{Episode} is
                    rolled out under the current policy with MDP depth budget
                    $\delta_{\max} = 12$; the resulting sketch $\pi$ is
                    filled with up to $N_{\mathrm{fill}} = 4$ random
                    substitutions via \texttt{\_fill\_in\_sketch}; the
                    substitution maximising $R$ is kept; the program is
                    inserted into the buffer and, if its fingerprint is novel,
                    appended to \texttt{rl\_sketches} and added to
                    \texttt{corpus\_fingerprints}.
              \item \emph{Training} (8 gradient steps of batch size 64): a
                    batch is drawn uniformly from the buffer and the
                    reward-weighted log-likelihood loss
                    $\mathcal{L}_{\mathrm{RL}}$ is minimised with Adam
                    ($\eta = 10^{-4}$) and gradient clipping to 1.0.
            \end{enumerate}
            Checkpoints of \texttt{rl\_sketches} are written every 10{,}000
            iterations.

      \item \emph{Post-RL expansion}.  The 100{,}000 novel RL
            sketches are expanded by \texttt{expand\_sketches} with
            \texttt{max\_substitutions} $= 100$ per sketch and
            \texttt{target\_total} $= 4{,}000{,}000$.  The expanded concrete
            programs are deduplicated against all fingerprints in
            \texttt{corpus\_fingerprints} (i.e.\ against the entire
            enumeration corpus and all RL-sampled programmes); only genuinely
            novel behaviours are retained.  This yields 3{,}906{,}215 novel
            programs, written to \texttt{rl\_corpus.json}.
    \end{itemize}

  \item \textbf{Final corpus.}  The two output files together form
    the full dataset:
    \[
      \underbrace{1{,}793{,}723}_{\texttt{enum\_corpus}}
        + \underbrace{3{,}906{,}215}_{\texttt{rl\_corpus}}
        = 5{,}699{,}938\;\text{programs.}
    \]

\end{itemize}

\paragraph{Program counts per stage.}

\begin{center}
\begin{tabular}{llrr}
\toprule
Stage & Output file & Sketches & Concrete programs \\
\midrule
Phase 1: Enumeration          & \texttt{enum\_sketches.json}  & 643{,}837 & --- \\
Phase 1: Quality filter       & (in-memory)                   & 581{,}984 & --- \\
Phase 2: Sketch expansion     & \texttt{enum\_corpus.json}    & --- & 1{,}793{,}723 \\
Phase 3: RL sampling          & (in-memory)                   & 100{,}000 & --- \\
Phase 3: Post-RL expansion    & \texttt{rl\_corpus.json}      & --- & 3{,}906{,}215 \\
\midrule
\textbf{Total}                &                               & & \textbf{5{,}699{,}938} \\
\bottomrule
\end{tabular}
\end{center}

\appendix

\section{Corpus Analysis}\label{sec:analysis}

This section characterises the generated corpus along three axes: structural
complexity, type and primitive coverage, and the diversity of the
enumeration and RL sub-corpora.

\subsection{Structural Complexity}\label{sec:complexity}

The two sub-corpora occupy largely disjoint regions of the program-size
spectrum.  Because the enumerator is bounded by $s_{\max} = 8$, its output is
concentrated at the upper end of that range: 83.9\% of enumeration programs
have size exactly~8, with a median of~8 and mean of~7.8.  The RL sub-corpus,
generated with MDP depth budget $\delta_{\max} = 12$, covers a much broader
range.  RL program sizes span 5--98, with a median of~22 and mean of~23.9.
The bulk of the RL corpus (43.2\%) falls in the size 21--35 range, while a
long tail extends to programs of size~98.  Table~\ref{tab:size-dist}
summarises the distribution.

\begin{table}[h]
\centering
\caption{Program size distribution across the two sub-corpora.}
\label{tab:size-dist}
\begin{tabular}{lrrrr}
\toprule
 & Min & Median & Mean & Max \\
\midrule
Enumeration ($n = 1{,}793{,}723$) & 1 & 8 & 7.8 & 8 \\
RL ($n = 3{,}906{,}215$)          & 5 & 22 & 23.9 & 98 \\
\bottomrule
\end{tabular}
\end{table}

The RL size distribution is approximately bell-shaped (Table~\ref{tab:rl-buckets}),
indicating that the policy learns to produce programs of moderate depth rather
than degenerately short or excessively long ones.

\begin{table}[h]
\centering
\caption{RL corpus breakdown by size bucket.}
\label{tab:rl-buckets}
\begin{tabular}{lrr}
\toprule
Size range & Programs & Fraction \\
\midrule
5--8   &     75{,}376 &  1.9\% \\
9--12  &    594{,}142 & 15.2\% \\
13--20 &    980{,}010 & 25.1\% \\
21--35 &  1{,}687{,}452 & 43.2\% \\
36--50 &    479{,}966 & 12.3\% \\
51--98 &     89{,}269 &  2.3\% \\
\bottomrule
\end{tabular}
\end{table}

\subsection{Type and Primitive Coverage}\label{sec:coverage}

Because the MDP generates programs of type
$\mathrm{list}[\mathrm{int}] \to \mathrm{list}[\mathrm{int}]$, the RL
sub-corpus is homogeneous in its return type.  Return-type diversity comes
entirely from the enumeration sub-corpus, which covers all five concrete
types: $\mathrm{list}[\mathrm{int}]$ (51.5\%),
$\mathrm{list}[\mathrm{list}[\mathrm{int}]]$ (23.0\%),
$\mathrm{int}$ (22.9\%), $\mathrm{list}[\mathrm{bool}]$ (2.5\%), and
$\mathrm{bool}$ (0.1\%).

Primitive coverage is near-complete in both sub-corpora: enumeration
programs use 48 of the grammar's primitives and RL programs use 47, with the
union covering all~48.  The two sub-corpora differ in \emph{how} primitives
are combined.  Higher-order functions are substantially more prevalent in the
RL corpus: lambda expressions appear 29$\times$ more frequently per
program in RL than in enumeration, and 3.3\% of RL programs
(128{,}414 programs) contain nested lambdas (i.e.\ reference inner
parameters \texttt{\_p1} or \texttt{\_p2}), compared to a negligible fraction
of enumeration programs.  Multi-argument list operations---\texttt{splice},
\texttt{insert}, \texttt{concat}, \texttt{cut\_slice},
\texttt{takelast}---are 5--8$\times$ more prevalent in the RL corpus, as
these operations require larger programs to use in compositionally
interesting ways.

\subsection{Diversity}\label{sec:diversity}

All programs have distinct fingerprints and all pass pass the quality filter (variability $\geq$ 0.3, non-crashing, non-constant). The enumeration corpus covers programs of size 1-8, while the RL helps extend beyond this size limit.

Some RL-generated programs suffer from internal redundancies. While their overall behaviour is interesting, they may be represented by shorter programs. For example, \texttt{product (repeat 10 8)} is an constant internal node in the AST. In the future, it may be worth simplifying RL-generated programs.

\bibliographystyle{plainnat}
\bibliography{references}

\end{document}
